%% Eigenständigkeit. Neu als Erklärung über den Einsatz von generativen KI-Modellem
% Das Zentrale Prüfungsamt fordert folgende Datei: https://www.rwth-aachen.de/global/show_document.asp?id=aaaaaaaaaaoasiu
\chapter*{Erklärung über den Einsatz von generativen KI-Modellen in der Arbeit} 

Hiermit versichere ich, dass ich KI-gestützte Werkzeuge zu Text- oder Bildgenerierung in der vorliegenden Arbeit ausschließlich in den in der folgenden Tabelle gekennzeichneten Bereichen zu den jeweils angegebenen Zwecken verwendet habe. Ausgenommen von dieser Erklärung ist der Einsatz von KI-gestützten Werkzeugen für Zwecke, die sich nicht auf den wissenschaftlichen Inhalt der Arbeit auswirken (bspw. Grammatik-, Stil-, Rechtschreibprüfung und  Übersetzungsprogramme).

\vspace*{1cm}
\begin{minipage}{\textwidth}
\centering
\begin{tabularx}{\textwidth-4em}{Xl}
\toprule	
Abschnitt der Arbeit & Einsatzzweck \\ 
\midrule	
\textit{Seitenzahlen/ Unterkapitel / Abbildungen}& \textit{Art der Bearbeitung}\\%Hier werden Seitenzahlen/ Unterkapitel / Abbildungen genannt, die mit KI-gestützten Tools erstellt/bearbeitet wurden und die Art der Bearbeitung genannt. 		
\bottomrule	
\end{tabularx}
\end{minipage}

\vspace*{1cm}
Mir ist bekannt, dass ich dazu verpflichtet bin, Ergebnisse von generativen KI-Modellen immer gründlich zu prüfen. Ich übernehme die Verantwortung für alle Ergebnisse, die ich in meine Arbeit einschließe. Weiterhin ist mir bewusst, dass die ungeprüfte Übernahme von Halluzinationen einer generativen KI in meine Arbeit als Täuschungsversuch interpretiert werden muss.

%I hereby assure that I composed the thesis at hand self-reliantly.
\selectlanguage{ngerman}
\vspace*{1cm}\begin{tabular}{p{7cm}p{5cm}}
Aachen, den \today & $\rule{5cm}{0.01cm}$ \\
& (Unterschrift)
\end{tabular}
